\label{gu:app:normalization}
\subsection{The phase and Fourier covariance}
We prove the covariance used in Proposition~\ref{gu:jac:extension}, with
the conventions of Section~\ref{gu:sec:setup}. An upper shear
$(x,y)\mapsto(x+Ay,y)$ acts by multiplication by
$\psi(\tfrac12[AT,T])$: conjugating \eqref{gu:eq:heisenberg-action} uses
\[
 [A(T+y),T+y]-[AT,T]=2[Ay,T]+[Ay,y].
\]
For the shear $n(u)$, the tensor pairing gives
\[
 [(1\otimes u)T,T]
 =-\tfrac12\Tr_{E/F}\tr_E(u\overline{G_T}).
\]
Its half-phase is \eqref{gu:jac:ordinary-n}, since the matrix trace lies
in $F$. The conjugate of $G_T$ transports the form from $X^c$ to $X$.

For the Weyl element
\[
 w=\begin{pmatrix}0&I_2\\-I_2&0\end{pmatrix},
\]
the operator is the Fourier transform
\begin{equation}\label{gu:jac:fourier}
 [\mathcal F_\beta f](T)
 =\int_Zf(S)\psi(\beta(T,S))\,dS,
 \qquad
 \beta(T,S)=-\frac12\operatorname{Tr}_{E/F}
                         \operatorname{tr}_E(T^\dagger J_4S),
\end{equation}
where $dS$ is self-dual for $\beta$ and $\psi$.  The partial Weyl elements
act by the corresponding partial Fourier transforms.  The Weil-index
coefficients in these formulas are $1$: the quadratic $F$-space underlying
the split Hermitian space $V$ is hyperbolic, a hyperbolic plane has Weil
index $1$, and the index is multiplicative under orthogonal direct sum.
The same assertion holds after multiplying the quadratic form or the
additive character by any nonzero scalar.

To prove \eqref{gu:jac:covariance}, the $U(V)$ formula is immediate;
left multiplication by $r_s$ commutes with the Levi's right action.
For the upper radical use
\[
 G_{r_s^{-1}T}=s^{-1}G_T,\qquad
 d_sn(u)d_s^{-1}=n(s^{-1}u).
\]
It remains to verify covariance at the Fourier operator.  The bilinear
form in \eqref{gu:jac:fourier} satisfies
\[
 \beta(r_s^{-1}T,r_s^{-1}S)=s^{-1}\beta(T,S).
\]
Using $U=r_sS$ in the Fourier integral therefore gives
\begin{align*}
 [D_s\mathcal F_\beta D_s^{-1}f](T)
 &=\int_Z f(r_sS)\psi(\beta(r_s^{-1}T,S))\,dS\\
 &=|s|_F^{-8}\int_Z f(U)\psi(s^{-1}\beta(T,U))\,dU\\
 &=|s|_F^{-8}[\mathcal F_\beta f](s^{-1}T).
\end{align*}
On the group side,
\[
 d_swd_s^{-1}=m(s^{-1}I_2,1)w,
 \qquad |\det(s^{-1}I_2)|_E^2=|s|_F^{-8}.
\]
The final integral is consequently the operator on the right
side of \eqref{gu:jac:covariance}. The Siegel parabolic and $w$ generate
both Siegel radicals; in each hyperbolic
coordinate plane, a product of upper, lower, and upper one-parameter
unipotents gives the partial Weyl element.  The Siegel Bruhat
decomposition then gives all of $U(W)$.  This proves
\eqref{gu:jac:covariance} on the whole isometry subgroup.

Equations \eqref{gu:jac:semidirect-decomposition} and \eqref{gu:jac:covariance},
together with $D_sD_t=D_{st}$, prove the representation law.
The operator $D_s$ is the canonical polarized Heisenberg intertwiner for
the symplectic transformation \eqref{gu:jac:tensor-action} associated with
$e_s$.  Thus \eqref{gu:jac:extension-definition} has Heisenberg
covariance.  Smoothness follows from the operator formulas above and the
smoothness of $\omega$.

\subsection{Schwartz restriction and the orbit quotient}\label{gu:app:topology}
\begin{lemma}\label{gu:jac:restriction}
Restriction from $M=F^\times\times Z$ to $C=\widetilde{\mathcal O}_B$
is surjective, its kernel is the algebraic $(N(X),\psi_B)$-relation
space, and the induced map is a homeomorphism for the Schwartz
inductive-limit topologies. In particular, the relation space is closed.
\end{lemma}
\begin{proof}
By \eqref{gu:jac:induced-n}, the unipotent character on $C$ is
$\psi_B$, so every relation vanishes there. If $f|_C=0$, local constancy makes the support
of $f$ a compact subset of $M\setminus C$. At every point of this support
some scalar character difference
\[
 d_n(c,T)=\psi\left(-\tfrac1{2c}\tr_E(u\overline{G_T})\right)
                  -\psi_B(n(u))
\]
is nonzero. A finite disjoint compact open refinement of these
neighborhoods writes $f=\sum_i f_i$, with $d_{n_i}$ nonzero on the
support of $f_i$. Extending $f_i/d_{n_i}$ by zero gives a Schwartz
function $g_i$ and
$f_i=(\Omega_0^{\GU}(n_i)-\psi_B(n_i))g_i$.
Conversely, a finite constancy partition of a Schwartz function on $C$
extends to compact open neighborhoods in $M$, chosen to avoid the other
values and the zero part of $C$. Disjoint refinement proves surjectivity.

The spaces under consideration are second-countable locally compact
totally disconnected spaces.  Their Schwartz spaces are inductive limits
of finite-dimensional spaces obtained by bounding the support and fixing
a finite compact open constancy partition.  Every finite-dimensional
subspace is contained in one such stage.  They therefore have the finest
locally convex topology on their underlying countable-dimensional vector
spaces.  Restriction is continuous, its kernel is closed, and the
algebraic quotient just constructed is a homeomorphism for these
topologies.  Equivalently, the extension argument supplies a continuous
lift on each finite-dimensional stage.
\end{proof}

\begin{lemma}\label{gu:jac:orbit-topology}
The transitive orbit map for \eqref{gu:jac:orbit-action} is a submersion.
Its induced bijection $L/H_j\to\widetilde{\mathcal O}_B$ is a
homeomorphism.
\end{lemma}
\begin{proof}
We verify the quotient topology. At a point of the scaled Gram locus,
$G_T=c\overline{\mathsf B}$ is invertible because $c\ne0$ and $B$ is
nondegenerate. The differential of the Gram map is
\[
 dG_T(S)=S^\dagger J_4T+T^\dagger J_4S.
\]
For any Hermitian matrix $D$, taking
$S=\tfrac12TG_T^{-1}D$ gives $dG_T(S)=D$.  Thus the equations defining
$\widetilde{\mathcal O}_B$ make it a smooth $F$-analytic manifold of
dimension $1+16-4=13$.

The orbit map from $G_\theta$ is a submersion.  It suffices to check this at
$(1,j)$.  A tangent vector $(\dot c,S)$ there satisfies
$dG_j(S)=\dot c\overline{\mathsf B}$. We use the decomposition
$V=j(X^c)\oplus j(X^c)^\perp$ and write $S=(b,C)$, identifying the first
summand by $j$.  With $*$ denoting adjoint for the forms on these two
nondegenerate summands, the tangent equation is $b^*+b=\dot c I$.
The block endomorphism
\[
 A=\begin{pmatrix}b&-C^*\\C&\dot c I/2\end{pmatrix}
\]
satisfies $A^*+A=\dot c I$ and $Aj=S$.  It is a tangent vector of $G_\theta$
with image $(\dot c,S)$.  Translation gives surjectivity of the
differential at every point.  The nonarchimedean inverse-function
theorem now gives local continuous analytic sections of the orbit map.
In particular, the orbit map is open.  Its induced bijection from the
quotient is therefore a homeomorphism.
\end{proof}
